\nivelbasico


\begin{ejercicio}
Un automóvil se mueve en línea recta. En cierto instante su velocidad es:

\[ v_i=\SI{4.0}{\meter\per\second}. \]

Seis segundos después su velocidad es:

\[ v_f=\SI{16.0}{\meter\per\second}. \]

Determine la aceleración media durante ese intervalo.

\textit{Pista: la aceleración media relaciona el cambio de velocidad con el intervalo de tiempo.}
\end{ejercicio}

\begin{respuesta}
\[ \boxed{\bar{a}_x=\SI{2.0}{\meter\per\second\squared}}. \]
\end{respuesta}

\begin{explicacion}
La aceleración media se define mediante:

\[ \bar{a}_x=\frac{\Delta v_x}{\Delta t}. \]

El cambio de velocidad es:

\[
\Delta v_x
=
v_f-v_i
=
\SI{16.0}{\meter\per\second}
-
\SI{4.0}{\meter\per\second}.
\]

Por tanto:

\[ \Delta v_x=\SI{12.0}{\meter\per\second}. \]

Como:

\[ \Delta t=\SI{6.0}{\second}, \]

tenemos:

\[
\bar{a}_x
=
\frac{\SI{12.0}{\meter\per\second}}
{\SI{6.0}{\second}}.
\]

Entonces:

\[ \boxed{\bar{a}_x=\SI{2.0}{\meter\per\second\squared}}. \]

Esto significa que, en promedio, la velocidad aumenta en:

\[ \SI{2.0}{\meter\per\second} \]

por cada segundo transcurrido.
\end{explicacion}


\begin{ejercicio}
Un automóvil se mueve en el sentido positivo del eje \(x\) con velocidad inicial:

\[ v_{0x}=\SI{20}{\meter\per\second}. \]

Después de \(\SI{4.0}{\second}\), su velocidad es:

\[ v_x=\SI{8.0}{\meter\per\second}. \]

Determine:

\begin{enumerate}[label=\alph*)]
    \item la aceleración media;
    \item si la rapidez aumenta o disminuye.
\end{enumerate}

\textit{Pista: compare los signos de la velocidad y de la aceleración.}
\end{ejercicio}

\begin{respuesta}
\[
\boxed{\bar{a}_x=-\SI{3.0}{\meter\per\second\squared}}.
\]

La rapidez disminuye.
\end{respuesta}

\begin{explicacion}
La aceleración media es:

\[
\bar{a}_x
=
\frac{v_x-v_{0x}}{\Delta t}.
\]

Sustituyendo:

\[
\bar{a}_x
=
\frac{
\SI{8.0}{\meter\per\second}
-
\SI{20}{\meter\per\second}
}{
\SI{4.0}{\second}
}.
\]

Entonces:

\[
\bar{a}_x
=
\frac{-\SI{12}{\meter\per\second}}
{\SI{4.0}{\second}}.
\]

Por tanto:

\[ \boxed{\bar{a}_x=-\SI{3.0}{\meter\per\second\squared}}. \]

La velocidad es positiva mientras que la aceleración es negativa.

Como tienen signos opuestos, la magnitud de la velocidad disminuye.

Por tanto:

\[ \boxed{\text{la rapidez disminuye}}. \]

Una aceleración negativa no significa por sí sola que el objeto esté frenando; en este caso la rapidez disminuye porque velocidad y aceleración tienen sentidos opuestos.
\end{explicacion}


\begin{ejercicio}
Una motocicleta parte del reposo y se mueve con aceleración constante:

\[ a_x=\SI{2.4}{\meter\per\second\squared}. \]

Determine, después de:

\[ t=\SI{5.0}{\second}, \]

\begin{enumerate}[label=\alph*)]
    \item su velocidad;
    \item su desplazamiento.
\end{enumerate}

\textit{Pista: partir del reposo significa \(v_{0x}=0\).}
\end{ejercicio}

\begin{respuesta}
\[
\boxed{v_x=\SI{12.0}{\meter\per\second}},
\qquad
\boxed{\Delta x=\SI{30.0}{\meter}}.
\]
\end{respuesta}

\begin{explicacion}
Los datos son:

\[
v_{0x}=0,
\qquad
a_x=\SI{2.4}{\meter\per\second\squared},
\qquad
t=\SI{5.0}{\second}.
\]

Para la velocidad utilizamos:

\[ v_x=v_{0x}+a_xt. \]

Sustituyendo:

\[
v_x
=
0+
(\SI{2.4}{\meter\per\second\squared})
(\SI{5.0}{\second}).
\]

Por tanto:

\[ \boxed{v_x=\SI{12.0}{\meter\per\second}}. \]

Para el desplazamiento:

\[ \Delta x=v_{0x}t+\frac{1}{2}a_xt^2. \]

Como \(v_{0x}=0\):

\[
\Delta x
=
\frac{1}{2}
(\SI{2.4}{\meter\per\second\squared})
(\SI{5.0}{\second})^2.
\]

Entonces:

\[ \boxed{\Delta x=\SI{30.0}{\meter}}. \]
\end{explicacion}


\begin{ejercicio}
Una partícula se encuentra inicialmente en:

\[ x_0=-\SI{4.0}{\meter} \]

y tiene velocidad inicial:

\[ v_{0x}=\SI{6.0}{\meter\per\second}. \]

Se mueve con aceleración constante:

\[ a_x=\SI{1.5}{\meter\per\second\squared}. \]

Determine después de \(\SI{8.0}{\second}\):

\begin{enumerate}[label=\alph*)]
    \item la velocidad;
    \item el desplazamiento;
    \item la posición.
\end{enumerate}

\textit{Pista: distinga entre desplazamiento \(\Delta x\) y posición final \(x\).}
\end{ejercicio}

\begin{respuesta}
\[
v_x=\SI{18.0}{\meter\per\second},
\qquad
\Delta x=\SI{96}{\meter},
\qquad
x=\SI{92}{\meter}.
\]
\end{respuesta}

\begin{explicacion}
Para la velocidad:

\[ v_x=v_{0x}+a_xt. \]

Sustituyendo:

\[
v_x
=
\SI{6.0}{\meter\per\second}
+
(\SI{1.5}{\meter\per\second\squared})
(\SI{8.0}{\second}).
\]

Por tanto:

\[ \boxed{v_x=\SI{18.0}{\meter\per\second}}. \]

El desplazamiento es:

\[ \Delta x=v_{0x}t+\frac{1}{2}a_xt^2. \]

Entonces:

\[
\Delta x
=
(6.0)(8.0)
+
\frac{1}{2}(1.5)(8.0)^2.
\]

En unidades SI:

\[
\Delta x
=
48+48
=
\SI{96}{\meter}.
\]

Así:

\[ \boxed{\Delta x=\SI{96}{\meter}}. \]

Finalmente:

\[ x=x_0+\Delta x. \]

Por tanto:

\[
x
=
-\SI{4.0}{\meter}
+
\SI{96}{\meter}.
\]

Entonces:

\[ \boxed{x=\SI{92}{\meter}}. \]
\end{explicacion}


\begin{ejercicio}
Un automóvil reduce uniformemente su velocidad desde:

\[ \SI{18}{\meter\per\second} \]

hasta:

\[ \SI{6.0}{\meter\per\second} \]

en un intervalo de:

\[ \SI{4.0}{\second}. \]

Determine:

\begin{enumerate}[label=\alph*)]
    \item la aceleración;
    \item el desplazamiento durante el intervalo.
\end{enumerate}

\textit{Pista: con aceleración constante puede utilizarse \(\bar{v}_x=(v_{0x}+v_x)/2\).}
\end{ejercicio}

\begin{respuesta}
\[
\boxed{a_x=-\SI{3.0}{\meter\per\second\squared}},
\qquad
\boxed{\Delta x=\SI{48}{\meter}}.
\]
\end{respuesta}

\begin{explicacion}
La aceleración es:

\[
a_x
=
\frac{v_x-v_{0x}}{t}.
\]

Sustituyendo:

\[
a_x
=
\frac{
\SI{6.0}{\meter\per\second}
-
\SI{18}{\meter\per\second}
}{
\SI{4.0}{\second}
}.
\]

Por tanto:

\[ \boxed{a_x=-\SI{3.0}{\meter\per\second\squared}}. \]

Como la aceleración es constante, la velocidad media durante el intervalo es:

\[
\bar{v}_x
=
\frac{v_{0x}+v_x}{2}.
\]

Entonces:

\[
\bar{v}_x
=
\frac{18+6}{2}
=
\SI{12}{\meter\per\second}.
\]

El desplazamiento es:

\[
\Delta x=\bar{v}_xt.
\]

Por tanto:

\[
\Delta x
=
(\SI{12}{\meter\per\second})
(\SI{4.0}{\second}).
\]

Así:

\[ \boxed{\Delta x=\SI{48}{\meter}}. \]

El desplazamiento continúa siendo positivo porque el automóvil no llegó a cambiar de sentido.
\end{explicacion}


\begin{ejercicio}
Una partícula se mueve con velocidad inicial:

\[ v_{0x}=\SI{5.0}{\meter\per\second} \]

y aceleración constante:

\[ a_x=\SI{2.0}{\meter\per\second\squared}. \]

Después de recorrer un desplazamiento:

\[ \Delta x=\SI{24}{\meter}, \]

¿qué velocidad tiene?

\textit{Pista: seleccione una ecuación de MRUA que no contenga el tiempo.}
\end{ejercicio}

\begin{respuesta}
\[ \boxed{v_x=\SI{11.0}{\meter\per\second}}. \]
\end{respuesta}

\begin{explicacion}
No se conoce el tiempo y tampoco es necesario calcularlo.

Utilizamos:

\[ v_x^2=v_{0x}^2+2a_x\Delta x. \]

Sustituyendo:

\[
v_x^2
=
(\SI{5.0}{\meter\per\second})^2
+
2
(\SI{2.0}{\meter\per\second\squared})
(\SI{24}{\meter}).
\]

Entonces:

\[
v_x^2
=
25+96
=
121\ \si{\meter\squared\per\second\squared}.
\]

Por tanto:

\[ v_x=\sqrt{121}\ \si{\meter\per\second}. \]

Así:

\[ \boxed{v_x=\SI{11.0}{\meter\per\second}}. \]

En este caso se toma la raíz positiva porque los datos describen un movimiento que continúa en el sentido positivo.
\end{explicacion}


\begin{ejercicio}
La gráfica velocidad--tiempo de una partícula es:

\begin{center}
\begin{tikzpicture}
\begin{axis}[
    width=0.78\textwidth,
    height=6cm,
    axis lines=left,
    xlabel={\(t\) (\si{\second})},
    ylabel={\(v_x\) (\si{\meter\per\second})},
    xmin=0,
    xmax=5,
    ymin=0,
    ymax=16,
    xtick={0,1,2,3,4,5},
    ytick={0,2,4,6,8,10,12,14,16},
    grid=both
]
\addplot[
    very thick,
    AzulPrincipal,
    domain=0:4,
    samples=2
]
{2+3*x};
\end{axis}
\end{tikzpicture}
\end{center}

Determine:

\begin{enumerate}[label=\alph*)]
    \item la velocidad inicial;
    \item la aceleración;
    \item la velocidad en \(t=\SI{4}{\second}\);
    \item el desplazamiento entre \(t=0\) y \(t=\SI{4}{\second}\).
\end{enumerate}

\textit{Pista: la pendiente de una gráfica \(v_x-t\) es la aceleración y el área bajo ella es el desplazamiento.}
\end{ejercicio}

\begin{respuesta}
\[
v_{0x}=\SI{2.0}{\meter\per\second},
\qquad
a_x=\SI{3.0}{\meter\per\second\squared},
\]

\[
v_x(\SI{4}{\second})=\SI{14}{\meter\per\second},
\qquad
\Delta x=\SI{32}{\meter}.
\]
\end{respuesta}

\begin{explicacion}
La gráfica comienza en:

\[ \boxed{v_{0x}=\SI{2.0}{\meter\per\second}}. \]

La aceleración es la pendiente:

\[
a_x
=
\frac{\Delta v_x}{\Delta t}.
\]

Tomando los puntos:

\[
(0,\SI{2}{\meter\per\second})
\]

y:

\[
(\SI{4}{\second},\SI{14}{\meter\per\second}),
\]

tenemos:

\[
a_x
=
\frac{14-2}{4}
=
\boxed{\SI{3.0}{\meter\per\second\squared}}.
\]

La velocidad final se lee directamente de la gráfica:

\[ \boxed{v_x=\SI{14}{\meter\per\second}}. \]

El desplazamiento corresponde al área bajo la gráfica.

Puede calcularse como el área de un trapecio:

\[
\Delta x
=
\frac{v_{0x}+v_x}{2}t.
\]

Entonces:

\[
\Delta x
=
\frac{2+14}{2}(4).
\]

Por tanto:

\[ \boxed{\Delta x=\SI{32}{\meter}}. \]
\end{explicacion}


\nivelintermedio


\begin{ejercicio}
Una partícula se encuentra inicialmente en:

\[ x_0=\SI{30}{\meter}. \]

Su velocidad inicial y aceleración son:

\[
v_{0x}=-\SI{4.0}{\meter\per\second},
\qquad
a_x=-\SI{2.5}{\meter\per\second\squared}.
\]

Determine después de \(\SI{6.0}{\second}\):

\begin{enumerate}[label=\alph*)]
    \item la velocidad;
    \item el desplazamiento;
    \item la posición final;
    \item si la rapidez aumentó o disminuyó.
\end{enumerate}

\textit{Pista: una aceleración negativa no significa necesariamente disminución de rapidez.}
\end{ejercicio}

\begin{respuesta}
\[
v_x=-\SI{19}{\meter\per\second},
\]

\[
\Delta x=-\SI{69}{\meter},
\qquad
x_f=-\SI{39}{\meter}.
\]

La rapidez aumenta.
\end{respuesta}

\begin{explicacion}
La velocidad final es:

\[
v_x=v_{0x}+a_xt.
\]

Sustituyendo:

\[
v_x
=
-\SI{4.0}{\meter\per\second}
+
(-\SI{2.5}{\meter\per\second\squared})
(\SI{6.0}{\second}).
\]

Entonces:

\[ \boxed{v_x=-\SI{19}{\meter\per\second}}. \]

El desplazamiento es:

\[
\Delta x
=
v_{0x}t+\frac{1}{2}a_xt^2.
\]

Por tanto:

\[
\Delta x
=
(-4.0)(6.0)
+
\frac{1}{2}(-2.5)(6.0)^2.
\]

Así:

\[
\Delta x
=
-24-45
=
\boxed{-\SI{69}{\meter}}.
\]

La posición final es:

\[
x_f=x_0+\Delta x.
\]

Entonces:

\[
x_f
=
\SI{30}{\meter}
-
\SI{69}{\meter}
=
\boxed{-\SI{39}{\meter}}.
\]

Inicialmente la rapidez era:

\[ \SI{4.0}{\meter\per\second}. \]

Finalmente es:

\[ \SI{19}{\meter\per\second}. \]

Como velocidad y aceleración son ambas negativas, poseen el mismo sentido y la rapidez aumenta.

\[
\boxed{\text{la rapidez aumenta}}.
\]
\end{explicacion}


\begin{ejercicio}
Un móvil parte desde:

\[ x_0=0 \]

con:

\[
v_{0x}=\SI{12}{\meter\per\second},
\qquad
a_x=-\SI{3.0}{\meter\per\second\squared}.
\]

Analice el movimiento durante los primeros \(\SI{6.0}{\second}\).

Determine:

\begin{enumerate}[label=\alph*)]
    \item el instante en que se detiene momentáneamente;
    \item la posición donde cambia de sentido;
    \item la velocidad en \(t=\SI{6.0}{\second}\);
    \item la posición en \(t=\SI{6.0}{\second}\);
    \item el desplazamiento durante los seis segundos;
    \item la distancia total recorrida durante ese intervalo.
\end{enumerate}

\textit{Pista: si existe un cambio de sentido, distancia y magnitud del desplazamiento dejan de coincidir.}
\end{ejercicio}

\begin{respuesta}
\[
t_{\mathrm{cambio}}=\SI{4.0}{\second},
\qquad
x_{\mathrm{cambio}}=\SI{24}{\meter},
\]

\[
v_x(\SI{6}{\second})=-\SI{6.0}{\meter\per\second},
\qquad
x(\SI{6}{\second})=\SI{18}{\meter},
\]

\[
\Delta x=\SI{18}{\meter},
\qquad
d=\SI{30}{\meter}.
\]
\end{respuesta}

\begin{explicacion}
El móvil cambia de sentido cuando su velocidad pasa por cero.

Utilizamos:

\[ v_x=v_{0x}+a_xt. \]

Entonces:

\[
0=12-3t.
\]

Por tanto:

\[ \boxed{t_{\mathrm{cambio}}=\SI{4.0}{\second}}. \]

La posición en ese instante es:

\[
x
=
v_{0x}t+\frac{1}{2}a_xt^2.
\]

Sustituyendo:

\[
x
=
12(4)
+
\frac{1}{2}(-3)(4)^2.
\]

Entonces:

\[
x
=
48-24
=
\boxed{\SI{24}{\meter}}.
\]

Ahora evaluamos el movimiento en \(t=\SI{6}{\second}\).

La velocidad es:

\[
v_x
=
12-3(6)
=
\boxed{-\SI{6.0}{\meter\per\second}}.
\]

La posición:

\[
x
=
12(6)
+
\frac{1}{2}(-3)(6)^2.
\]

Entonces:

\[
x
=
72-54
=
\boxed{\SI{18}{\meter}}.
\]

Como comenzó en el origen:

\[ \boxed{\Delta x=\SI{18}{\meter}}. \]

Sin embargo, para calcular la distancia debemos considerar el cambio de sentido.

Primero recorre:

\[ \SI{24}{\meter} \]

hasta el punto de retorno.

Después regresa desde \(x=24\) hasta \(x=18\):

\[ \SI{6}{\meter}. \]

Así:

\[
d
=
24+6
=
\boxed{\SI{30}{\meter}}.
\]

Este ejercicio muestra por qué es indispensable comprobar si la velocidad cambia de signo dentro del intervalo.
\end{explicacion}


\begin{ejercicio}
La posición de una partícula está dada por:

\[ x(t)=2.0+6.0t-2.0t^2, \]

donde \(x\) se expresa en metros y \(t\) en segundos.

Determine:

\begin{enumerate}[label=\alph*)]
    \item la función velocidad;
    \item la aceleración;
    \item el instante en que cambia de sentido;
    \item la posición máxima;
    \item la posición y la velocidad en \(t=\SI{3.0}{\second}\);
    \item el desplazamiento y la distancia recorrida entre \(t=0\) y \(t=\SI{3.0}{\second}\).
\end{enumerate}

\textit{Pista: obtenga primero \(v_x(t)=\deriv{x}{t}\) y \(a_x=\deriv{v_x}{t}\).}
\end{ejercicio}

\begin{respuesta}
\[
v_x(t)=6.0-4.0t,
\qquad
a_x=-\SI{4.0}{\meter\per\second\squared},
\]

\[
t_{\mathrm{cambio}}=\SI{1.50}{\second},
\qquad
x_{\max}=\SI{6.50}{\meter},
\]

\[
x(\SI{3}{\second})=\SI{2.0}{\meter},
\qquad
v_x(\SI{3}{\second})=-\SI{6.0}{\meter\per\second},
\]

\[
\Delta x=0,
\qquad
d=\SI{9.0}{\meter}.
\]
\end{respuesta}

\begin{explicacion}
La velocidad instantánea es:

\[ v_x=\deriv{x}{t}. \]

Derivando:

\[
v_x(t)
=
6.0-4.0t.
\]

Por tanto:

\[ \boxed{v_x(t)=6.0-4.0t}. \]

La aceleración es:

\[ a_x=\deriv{v_x}{t}. \]

Entonces:

\[ \boxed{a_x=-\SI{4.0}{\meter\per\second\squared}}. \]

El cambio de sentido ocurre cuando:

\[ v_x=0. \]

Por tanto:

\[
0=6.0-4.0t.
\]

De donde:

\[ \boxed{t=\SI{1.50}{\second}}. \]

Evaluamos la posición:

\[
x(1.5)
=
2+6(1.5)-2(1.5)^2.
\]

Entonces:

\[
x_{\max}
=
2+9-4.5
=
\boxed{\SI{6.50}{\meter}}.
\]

En \(t=\SI{3}{\second}\):

\[
x(3)
=
2+18-18
=
\boxed{\SI{2.0}{\meter}}.
\]

La velocidad es:

\[
v_x(3)
=
6-4(3)
=
\boxed{-\SI{6.0}{\meter\per\second}}.
\]

Como:

\[
x(0)=\SI{2.0}{\meter}
\]

y:

\[
x(3)=\SI{2.0}{\meter},
\]

el desplazamiento es:

\[ \boxed{\Delta x=0}. \]

Para la distancia, la partícula se desplaza desde:

\[
x=\SI{2.0}{\meter}
\]

hasta:

\[
x=\SI{6.50}{\meter},
\]

recorriendo:

\[ \SI{4.50}{\meter}. \]

Luego regresa a:

\[ x=\SI{2.0}{\meter}, \]

recorriendo otros:

\[ \SI{4.50}{\meter}. \]

Por tanto:

\[ \boxed{d=\SI{9.0}{\meter}}. \]
\end{explicacion}


\begin{ejercicio}
Un automóvil aumenta uniformemente su rapidez desde:

\[ \SI{72}{\kilo\meter\per\hour} \]

hasta:

\[ \SI{108}{\kilo\meter\per\hour} \]

en:

\[ \SI{5.0}{\second}. \]

Determine:

\begin{enumerate}[label=\alph*)]
    \item las velocidades inicial y final en unidades SI;
    \item la aceleración;
    \item el desplazamiento durante el intervalo.
\end{enumerate}

\textit{Pista: antes de utilizar las ecuaciones cinemáticas convierta las velocidades a \(\si{\meter\per\second}\).}
\end{ejercicio}

\begin{respuesta}
\[
v_{0x}=\SI{20}{\meter\per\second},
\qquad
v_x=\SI{30}{\meter\per\second},
\]

\[
a_x=\SI{2.0}{\meter\per\second\squared},
\qquad
\Delta x=\SI{125}{\meter}.
\]
\end{respuesta}

\begin{explicacion}
Utilizamos:

\[
\SI{1}{\kilo\meter\per\hour}
=
\frac{1}{3.6}\si{\meter\per\second}.
\]

Por tanto:

\[
\SI{72}{\kilo\meter\per\hour}
=
\frac{72}{3.6}
=
\boxed{\SI{20}{\meter\per\second}}.
\]

Asimismo:

\[
\SI{108}{\kilo\meter\per\hour}
=
\frac{108}{3.6}
=
\boxed{\SI{30}{\meter\per\second}}.
\]

La aceleración es:

\[
a_x
=
\frac{v_x-v_{0x}}{t}.
\]

Entonces:

\[
a_x
=
\frac{30-20}{5.0}
=
\boxed{\SI{2.0}{\meter\per\second\squared}}.
\]

Como la aceleración es constante:

\[
\bar{v}_x
=
\frac{20+30}{2}
=
\SI{25}{\meter\per\second}.
\]

El desplazamiento es:

\[
\Delta x
=
\bar{v}_xt.
\]

Entonces:

\[
\Delta x
=
(\SI{25}{\meter\per\second})
(\SI{5.0}{\second}).
\]

Por tanto:

\[ \boxed{\Delta x=\SI{125}{\meter}}. \]
\end{explicacion}


\begin{ejercicio}
Una partícula se mueve con aceleración constante. Se sabe que:

\[
v_x(\SI{2.0}{\second})=\SI{11}{\meter\per\second},
\]

y:

\[
v_x(\SI{6.0}{\second})=\SI{23}{\meter\per\second}.
\]

Además:

\[ x_0=-\SI{10}{\meter}. \]

Determine:

\begin{enumerate}[label=\alph*)]
    \item la aceleración;
    \item la velocidad inicial;
    \item el desplazamiento entre \(t=0\) y \(t=\SI{6.0}{\second}\);
    \item la posición en \(t=\SI{6.0}{\second}\).
\end{enumerate}

\textit{Pista: la pendiente de una relación lineal \(v_x-t\) permite encontrar primero la aceleración.}
\end{ejercicio}

\begin{respuesta}
\[
a_x=\SI{3.0}{\meter\per\second\squared},
\qquad
v_{0x}=\SI{5.0}{\meter\per\second},
\]

\[
\Delta x=\SI{84}{\meter},
\qquad
x(\SI{6}{\second})=\SI{74}{\meter}.
\]
\end{respuesta}

\begin{explicacion}
Con aceleración constante:

\[
a_x
=
\frac{\Delta v_x}{\Delta t}.
\]

Por tanto:

\[
a_x
=
\frac{
\SI{23}{\meter\per\second}
-
\SI{11}{\meter\per\second}
}{
\SI{6}{\second}
-
\SI{2}{\second}
}.
\]

Entonces:

\[ \boxed{a_x=\SI{3.0}{\meter\per\second\squared}}. \]

Ahora usamos:

\[ v_x=v_{0x}+a_xt. \]

En \(t=\SI{2}{\second}\):

\[
11=v_{0x}+3(2).
\]

De donde:

\[ \boxed{v_{0x}=\SI{5.0}{\meter\per\second}}. \]

El desplazamiento durante los primeros seis segundos es:

\[
\Delta x
=
v_{0x}t+\frac{1}{2}a_xt^2.
\]

Entonces:

\[
\Delta x
=
5(6)
+
\frac{1}{2}(3)(6)^2.
\]

Así:

\[
\Delta x
=
30+54
=
\boxed{\SI{84}{\meter}}.
\]

La posición final es:

\[
x=x_0+\Delta x.
\]

Por tanto:

\[
x
=
-\SI{10}{\meter}
+
\SI{84}{\meter}.
\]

Entonces:

\[ \boxed{x=\SI{74}{\meter}}. \]
\end{explicacion}


\begin{ejercicio}
La velocidad de una partícula varía de acuerdo con la siguiente gráfica:

\begin{center}
\begin{tikzpicture}
\begin{axis}[
    width=0.82\textwidth,
    height=6.5cm,
    axis lines=middle,
    xlabel={\(t\) (\si{\second})},
    ylabel={\(v_x\) (\si{\meter\per\second})},
    xmin=0,
    xmax=5.5,
    ymin=-10,
    ymax=14,
    xtick={0,1,2,3,4,5},
    ytick={-8,-4,0,4,8,12},
    grid=both
]

\addplot[
    very thick,
    AzulPrincipal,
    domain=0:5,
    samples=2
]
{12-4*x};

\end{axis}
\end{tikzpicture}
\end{center}

Determine:

\begin{enumerate}[label=\alph*)]
    \item la aceleración;
    \item el instante en que cambia de sentido;
    \item la velocidad en \(t=\SI{5}{\second}\);
    \item el desplazamiento total entre \(0\) y \(\SI{5}{\second}\);
    \item la distancia total recorrida.
\end{enumerate}

\textit{Pista: cuando la gráfica cruza el eje temporal, la velocidad cambia de signo.}
\end{ejercicio}

\begin{respuesta}
\[
a_x=-\SI{4.0}{\meter\per\second\squared},
\qquad
t_{\mathrm{cambio}}=\SI{3.0}{\second},
\]

\[
v_x(\SI{5}{\second})=-\SI{8.0}{\meter\per\second},
\]

\[
\Delta x=\SI{10}{\meter},
\qquad
d=\SI{26}{\meter}.
\]
\end{respuesta}

\begin{explicacion}
La aceleración es la pendiente de la gráfica:

\[
a_x
=
\frac{-8-12}{5-0}.
\]

Entonces:

\[ \boxed{a_x=-\SI{4.0}{\meter\per\second\squared}}. \]

El cambio de sentido ocurre cuando:

\[ v_x=0. \]

De la ecuación de la recta:

\[ v_x=12-4t. \]

Por tanto:

\[
0=12-4t,
\]

de donde:

\[ \boxed{t=\SI{3.0}{\second}}. \]

En \(t=\SI{5}{\second}\):

\[
v_x=12-4(5)=-8.
\]

Así:

\[ \boxed{v_x=-\SI{8.0}{\meter\per\second}}. \]

Para el desplazamiento total podemos utilizar la velocidad media, pues la aceleración es constante:

\[
\bar{v}_x
=
\frac{v_{0x}+v_x}{2}.
\]

Entonces:

\[
\bar{v}_x
=
\frac{12+(-8)}{2}
=
\SI{2.0}{\meter\per\second}.
\]

Por tanto:

\[
\Delta x
=
(\SI{2.0}{\meter\per\second})
(\SI{5.0}{\second})
=
\boxed{\SI{10}{\meter}}.
\]

Para la distancia debemos separar los dos sentidos de movimiento.

Durante los primeros \(\SI{3}{\second}\), el área bajo la gráfica es un triángulo:

\[
d_1
=
\frac{1}{2}(3)(12)
=
\SI{18}{\meter}.
\]

Entre \(t=3\) y \(t=5\), el área se encuentra debajo del eje. Su magnitud es:

\[
d_2
=
\frac{1}{2}(2)(8)
=
\SI{8}{\meter}.
\]

Por tanto:

\[
d
=
18+8
=
\boxed{\SI{26}{\meter}}.
\]

El área algebraica produce desplazamiento; la suma de las magnitudes de las áreas produce distancia.
\end{explicacion}


\begin{ejercicio}
Un móvil parte del reposo desde:

\[ x_0=-\SI{5}{\meter}. \]

Durante los primeros \(\SI{4.0}{\second}\) tiene aceleración constante:

\[ a_1=\SI{2.0}{\meter\per\second\squared}. \]

A continuación mantiene la velocidad alcanzada, pero comienza a acelerar en sentido contrario con:

\[ a_2=-\SI{1.0}{\meter\per\second\squared} \]

durante otros \(\SI{6.0}{\second}\).

Determine:

\begin{enumerate}[label=\alph*)]
    \item la velocidad y posición al finalizar la primera etapa;
    \item la velocidad y posición al finalizar la segunda;
    \item el desplazamiento total;
    \item la velocidad media durante los \(\SI{10}{\second}\).
\end{enumerate}

\textit{Pista: la velocidad final de la primera etapa es la velocidad inicial de la segunda.}
\end{ejercicio}

\begin{respuesta}
Primera etapa:

\[
v_1=\SI{8.0}{\meter\per\second},
\qquad
x_1=\SI{11}{\meter}.
\]

Segunda etapa:

\[
v_2=\SI{2.0}{\meter\per\second},
\qquad
x_2=\SI{41}{\meter}.
\]

Además:

\[
\Delta x=\SI{46}{\meter},
\qquad
\bar{v}_x=\SI{4.6}{\meter\per\second}.
\]
\end{respuesta}

\begin{explicacion}
En la primera etapa:

\[
v_{0x}=0,
\qquad
a_1=\SI{2.0}{\meter\per\second\squared},
\qquad
t_1=\SI{4.0}{\second}.
\]

La velocidad es:

\[
v_1
=
0+2(4)
=
\boxed{\SI{8.0}{\meter\per\second}}.
\]

El desplazamiento:

\[
\Delta x_1
=
\frac{1}{2}(2)(4)^2
=
\SI{16}{\meter}.
\]

Como:

\[ x_0=-\SI{5}{\meter}, \]

la posición es:

\[
x_1
=
-5+16
=
\boxed{\SI{11}{\meter}}.
\]

En la segunda etapa:

\[
v_{0,2}=\SI{8.0}{\meter\per\second},
\qquad
a_2=-\SI{1.0}{\meter\per\second\squared},
\qquad
t_2=\SI{6.0}{\second}.
\]

La velocidad final es:

\[
v_2
=
8-1(6)
=
\boxed{\SI{2.0}{\meter\per\second}}.
\]

El desplazamiento de esta etapa:

\[
\Delta x_2
=
8(6)
+
\frac{1}{2}(-1)(6)^2.
\]

Entonces:

\[
\Delta x_2
=
48-18
=
\SI{30}{\meter}.
\]

La posición final es:

\[
x_2
=
11+30
=
\boxed{\SI{41}{\meter}}.
\]

El desplazamiento total es:

\[
\Delta x
=
41-(-5)
=
\boxed{\SI{46}{\meter}}.
\]

El tiempo total es:

\[ \Delta t=\SI{10}{\second}. \]

Por tanto:

\[
\bar{v}_x
=
\frac{\SI{46}{\meter}}
{\SI{10}{\second}}
=
\boxed{\SI{4.6}{\meter\per\second}}.
\]
\end{explicacion}


\nivelavanzado


\begin{ejercicio}
Un automóvil se mueve inicialmente con velocidad \(v_0>0\). Al aplicar los frenos experimenta una aceleración constante:

\[ a_x=-b, \]

donde \(b>0\).

Deduzca expresiones para:

\begin{enumerate}[label=\alph*)]
    \item el tiempo necesario para detenerse;
    \item la distancia de frenado;
    \item la razón entre las distancias de frenado de dos automóviles idénticos cuyas velocidades iniciales son \(v_0\) y \(2v_0\).
\end{enumerate}

Interprete físicamente el último resultado.
\end{ejercicio}

\begin{respuesta}
\[
\boxed{t_{\mathrm{f}}=\frac{v_0}{b}},
\]

\[
\boxed{d_{\mathrm{f}}=\frac{v_0^2}{2b}},
\]

y:

\[
\boxed{\frac{d_{\mathrm{f},2}}{d_{\mathrm{f},1}}=4}.
\]

Duplicar la velocidad inicial cuadruplica la distancia de frenado, si \(b\) permanece constante.
\end{respuesta}

\begin{explicacion}
El automóvil se detiene cuando:

\[ v_x=0. \]

La ecuación de velocidad es:

\[ v_x=v_0+a_xt. \]

Como:

\[ a_x=-b, \]

tenemos:

\[ 0=v_0-bt_{\mathrm{f}}. \]

Por tanto:

\[ \boxed{t_{\mathrm{f}}=\frac{v_0}{b}}. \]

Para obtener la distancia de frenado utilizamos:

\[ v_x^2=v_0^2+2a_x\Delta x. \]

En la detención:

\[ v_x=0. \]

Entonces:

\[
0=v_0^2-2bd_{\mathrm{f}}.
\]

Despejando:

\[ \boxed{d_{\mathrm{f}}=\frac{v_0^2}{2b}}. \]

Para el segundo automóvil:

\[ v_{0,2}=2v_0. \]

Entonces:

\[
d_{\mathrm{f},2}
=
\frac{(2v_0)^2}{2b}
=
\frac{4v_0^2}{2b}.
\]

Mientras que:

\[
d_{\mathrm{f},1}
=
\frac{v_0^2}{2b}.
\]

Por tanto:

\[
\frac{d_{\mathrm{f},2}}{d_{\mathrm{f},1}}
=
4.
\]

Así:

\[ \boxed{d_{\mathrm{f},2}=4d_{\mathrm{f},1}}. \]

La distancia de frenado no es proporcional directamente a la velocidad inicial, sino a su cuadrado.

Por ello, un aumento moderado de velocidad puede producir un incremento considerable en la distancia necesaria para detener un vehículo.
\end{explicacion}


\begin{ejercicio}
Un automóvil circula por una carretera recta con rapidez:

\[ v_0=\SI{25.0}{\meter\per\second}. \]

El conductor observa un obstáculo. Su tiempo de reacción es:

\[ t_r=\SI{0.75}{\second}. \]

Durante ese tiempo mantiene aproximadamente su velocidad constante. Al comenzar a frenar, el automóvil experimenta una aceleración constante:

\[ a_x=-\SI{6.0}{\meter\per\second\squared}. \]

Determine:

\begin{enumerate}[label=\alph*)]
    \item la distancia recorrida durante el tiempo de reacción;
    \item el tiempo de frenado;
    \item la distancia recorrida mientras frena;
    \item el tiempo total hasta detenerse;
    \item la distancia total necesaria para detenerse desde que el conductor observa el obstáculo.
\end{enumerate}
\end{ejercicio}

\begin{respuesta}
\[
d_r=\SI{18.75}{\meter},
\]

\[
t_f\approx\SI{4.17}{\second},
\qquad
d_f\approx\SI{52.1}{\meter},
\]

\[
t_{\mathrm{total}}\approx\SI{4.92}{\second},
\qquad
d_{\mathrm{total}}\approx\SI{70.8}{\meter}.
\]
\end{respuesta}

\begin{explicacion}
El movimiento debe dividirse en dos etapas físicamente diferentes.

\subsubsection*{Etapa 1: tiempo de reacción}

Durante el tiempo de reacción se supone velocidad aproximadamente constante:

\[
d_r=v_0t_r.
\]

Entonces:

\[
d_r
=
(\SI{25.0}{\meter\per\second})
(\SI{0.75}{\second}).
\]

Por tanto:

\[ \boxed{d_r=\SI{18.75}{\meter}}. \]

\subsubsection*{Etapa 2: frenado}

Al comenzar a frenar:

\[
v_{0x}=\SI{25.0}{\meter\per\second},
\qquad
a_x=-\SI{6.0}{\meter\per\second\squared}.
\]

La detención ocurre cuando:

\[ v_x=0. \]

Entonces:

\[
0
=
25.0-6.0t_f.
\]

Por tanto:

\[
\boxed{t_f\approx\SI{4.17}{\second}}.
\]

Para la distancia de frenado utilizamos:

\[
v_x^2=v_{0x}^2+2a_xd_f.
\]

Así:

\[
0
=
(25.0)^2
+
2(-6.0)d_f.
\]

Entonces:

\[
d_f
=
\frac{625}{12}
\approx
\boxed{\SI{52.1}{\meter}}.
\]

El tiempo total es:

\[
t_{\mathrm{total}}
=
t_r+t_f.
\]

Por tanto:

\[
t_{\mathrm{total}}
=
0.75+4.17
\approx
\boxed{\SI{4.92}{\second}}.
\]

La distancia total es:

\[
d_{\mathrm{total}}
=
d_r+d_f.
\]

Así:

\[
d_{\mathrm{total}}
=
18.75+52.1
\approx
\boxed{\SI{70.8}{\meter}}.
\]

La distancia total de detención no es únicamente la distancia de frenado. El vehículo también avanza durante el tiempo de reacción del conductor.
\end{explicacion}


\begin{ejercicio}
En un experimento se registran las posiciones de una partícula en intervalos iguales de \(\SI{1.0}{\second}\):

\begin{center}
\begin{tabular}{c|ccccc}
\toprule
\(t\) (\si{\second}) & 0 & 1 & 2 & 3 & 4 \\
\midrule
\(x\) (\si{\meter}) & 0 & 3 & 10 & 21 & ? \\
\bottomrule
\end{tabular}
\end{center}

Suponiendo que el movimiento tiene aceleración constante:

\begin{enumerate}[label=\alph*)]
    \item utilice las diferencias entre posiciones sucesivas para justificar que los primeros cuatro datos son compatibles con MRUA;
    \item determine la aceleración;
    \item determine la velocidad inicial;
    \item prediga la posición en \(t=\SI{4.0}{\second}\);
    \item determine la velocidad en ese instante.
\end{enumerate}
\end{ejercicio}

\begin{respuesta}
Las diferencias de posición son:

\[
3,\quad 7,\quad 11\ \si{\meter},
\]

y sus diferencias sucesivas son constantes:

\[
4,\quad4\ \si{\meter}.
\]

Se obtiene:

\[
a_x=\SI{4.0}{\meter\per\second\squared},
\qquad
v_{0x}=\SI{1.0}{\meter\per\second},
\]

\[
x(\SI{4}{\second})=\SI{36}{\meter},
\qquad
v_x(\SI{4}{\second})=\SI{17}{\meter\per\second}.
\]
\end{respuesta}

\begin{explicacion}
Primero calculamos los desplazamientos durante cada segundo:

\[
\Delta x_1=3-0=\SI{3}{\meter},
\]

\[
\Delta x_2=10-3=\SI{7}{\meter},
\]

\[
\Delta x_3=21-10=\SI{11}{\meter}.
\]

Ahora calculamos cómo cambian esos desplazamientos:

\[
7-3=\SI{4}{\meter},
\]

\[
11-7=\SI{4}{\meter}.
\]

La segunda diferencia es constante.

Para intervalos temporales iguales \(\Delta t=\SI{1}{\second}\), en un movimiento con aceleración constante la segunda diferencia de posición satisface:

\[
\Delta^2x=a_x(\Delta t)^2.
\]

Entonces:

\[
a_x
=
\frac{\SI{4}{\meter}}
{(\SI{1}{\second})^2}.
\]

Por tanto:

\[ \boxed{a_x=\SI{4.0}{\meter\per\second\squared}}. \]

La ecuación de posición es:

\[
x=x_0+v_{0x}t+\frac{1}{2}a_xt^2.
\]

Como:

\[ x_0=0, \]

utilizamos el dato correspondiente a \(t=\SI{1}{\second}\):

\[
3
=
v_{0x}(1)
+
\frac{1}{2}(4)(1)^2.
\]

Entonces:

\[
3=v_{0x}+2.
\]

Por tanto:

\[ \boxed{v_{0x}=\SI{1.0}{\meter\per\second}}. \]

Para \(t=\SI{4}{\second}\):

\[
x
=
1(4)
+
\frac{1}{2}(4)(4)^2.
\]

Entonces:

\[
x
=
4+32
=
\boxed{\SI{36}{\meter}}.
\]

La velocidad es:

\[
v_x=v_{0x}+a_xt.
\]

Por tanto:

\[
v_x
=
1+4(4)
=
\boxed{\SI{17}{\meter\per\second}}.
\]

El ejercicio muestra que una tabla de posiciones puede utilizarse para reconocer indirectamente una aceleración constante.
\end{explicacion}


\begin{ejercicio}
Un automóvil pasa frente a una patrulla estacionada con velocidad constante:

\[ v_A=\SI{12.0}{\meter\per\second}. \]

En el instante en que el automóvil pasa frente a ella se define:

\[ t=0,\qquad x=0. \]

La patrulla permanece en reposo durante \(\SI{3.0}{\second}\). Después inicia la persecución desde el reposo con aceleración constante:

\[ a_P=\SI{4.0}{\meter\per\second\squared}. \]

Suponga que el automóvil mantiene su velocidad constante.

Determine:

\begin{enumerate}[label=\alph*)]
    \item cuánto tiempo después de que el automóvil pasa frente a la patrulla ocurre el alcance;
    \item cuánto tiempo dura realmente la persecución;
    \item la posición donde ocurre el alcance;
    \item la velocidad de la patrulla en ese instante.
\end{enumerate}
\end{ejercicio}

\begin{respuesta}
\[
t_{\mathrm{enc}}\approx\SI{11.2}{\second},
\]

\[
t_{\mathrm{persec}}\approx\SI{8.20}{\second},
\]

\[
x_{\mathrm{enc}}\approx\SI{134}{\meter},
\]

\[
v_P\approx\SI{32.8}{\meter\per\second}.
\]
\end{respuesta}

\begin{explicacion}
Tomaremos \(t\) desde el instante en que el automóvil pasa frente a la patrulla.

La posición del automóvil es:

\[
x_A=12t.
\]

La patrulla comienza a moverse sólo cuando:

\[ t=\SI{3.0}{\second}. \]

Por tanto, el tiempo durante el cual la patrulla ha estado acelerando es:

\[ t_P=t-3. \]

Como parte del reposo:

\[
x_P
=
\frac{1}{2}a_P(t-3)^2.
\]

Entonces:

\[
x_P
=
2(t-3)^2.
\]

En el instante del alcance:

\[ x_A=x_P. \]

Por tanto:

\[
12t=2(t-3)^2.
\]

Dividiendo entre 2:

\[
6t=(t-3)^2.
\]

Desarrollando:

\[
6t=t^2-6t+9.
\]

Entonces:

\[
t^2-12t+9=0.
\]

Aplicando la fórmula cuadrática:

\[
t
=
\frac{12\pm\sqrt{144-36}}{2}.
\]

Así:

\[
t=6\pm3\sqrt{3}.
\]

Se obtienen aproximadamente:

\[
t_1=\SI{0.804}{\second},
\]

y:

\[
t_2=\SI{11.196}{\second}.
\]

La primera raíz no es físicamente válida para la ecuación utilizada, pues corresponde a un instante anterior al inicio de la persecución:

\[ t<\SI{3.0}{\second}. \]

Por tanto:

\[ \boxed{t_{\mathrm{enc}}\approx\SI{11.2}{\second}}. \]

La patrulla estuvo realmente en movimiento durante:

\[
t_P
=
11.196-3.0.
\]

Entonces:

\[ \boxed{t_P\approx\SI{8.20}{\second}}. \]

La posición de encuentro puede obtenerse con la ecuación del automóvil:

\[
x_A
=
(\SI{12.0}{\meter\per\second})
(\SI{11.196}{\second}).
\]

Así:

\[ \boxed{x_{\mathrm{enc}}\approx\SI{134}{\meter}}. \]

La velocidad de la patrulla es:

\[
v_P
=
a_Pt_P.
\]

Entonces:

\[
v_P
=
(\SI{4.0}{\meter\per\second\squared})
(\SI{8.196}{\second}).
\]

Por tanto:

\[ \boxed{v_P\approx\SI{32.8}{\meter\per\second}}. \]

Este problema exige distinguir cuidadosamente entre el tiempo medido desde que pasa el automóvil y el tiempo durante el cual la patrulla realmente acelera.
\end{explicacion}


\begin{ejercicio}
La posición de una partícula está descrita por:

\[ x(t)=10t-t^2, \]

donde \(x\) está en metros y \(t\) en segundos.

Determine:

\begin{enumerate}[label=\alph*)]
    \item los dos instantes en que la partícula se encuentra en \(x=\SI{16}{\meter}\);
    \item la velocidad en cada uno de esos instantes;
    \item el instante y la posición donde cambia de sentido;
    \item explique físicamente por qué una misma posición puede corresponder a dos velocidades diferentes.
\end{enumerate}
\end{ejercicio}

\begin{respuesta}
\[
t_1=\SI{2.0}{\second},
\qquad
t_2=\SI{8.0}{\second},
\]

\[
v_x(t_1)=+\SI{6.0}{\meter\per\second},
\qquad
v_x(t_2)=-\SI{6.0}{\meter\per\second},
\]

\[
t_{\mathrm{cambio}}=\SI{5.0}{\second},
\qquad
x_{\max}=\SI{25}{\meter}.
\]
\end{respuesta}

\begin{explicacion}
Buscamos los instantes en que:

\[ x=\SI{16}{\meter}. \]

Entonces:

\[
16=10t-t^2.
\]

Reordenando:

\[
t^2-10t+16=0.
\]

Factorizamos:

\[
(t-2)(t-8)=0.
\]

Por tanto:

\[
\boxed{t_1=\SI{2.0}{\second}},
\qquad
\boxed{t_2=\SI{8.0}{\second}}.
\]

Ahora obtenemos la velocidad:

\[
v_x=\deriv{x}{t}.
\]

Por tanto:

\[
v_x=10-2t.
\]

En \(t=\SI{2}{\second}\):

\[
v_x=10-4
=
\boxed{+\SI{6.0}{\meter\per\second}}.
\]

En \(t=\SI{8}{\second}\):

\[
v_x=10-16
=
\boxed{-\SI{6.0}{\meter\per\second}}.
\]

El cambio de sentido ocurre cuando:

\[ v_x=0. \]

Así:

\[
0=10-2t.
\]

Entonces:

\[ \boxed{t_{\mathrm{cambio}}=\SI{5.0}{\second}}. \]

La posición es:

\[
x(5)
=
10(5)-(5)^2.
\]

Por tanto:

\[ \boxed{x_{\max}=\SI{25}{\meter}}. \]

La partícula atraviesa \(x=\SI{16}{\meter}\) una primera vez mientras se mueve hacia el sentido positivo y una segunda vez después de haber cambiado de sentido.

La posición no determina por sí sola el estado completo de movimiento. Para ello también necesitamos conocer la velocidad.
\end{explicacion}


\begin{ejercicio}
Partiendo únicamente de las relaciones válidas para aceleración constante:

\[
v_x=v_{0x}+a_xt
\]

y:

\[
\Delta x=\frac{v_{0x}+v_x}{2}\,t,
\]

deduzca la ecuación:

\[ v_x^2=v_{0x}^2+2a_x\Delta x. \]

Explique además por qué esta relación resulta especialmente útil en ciertos problemas de MRUA.
\end{ejercicio}

\begin{respuesta}
La deducción conduce a:

\[ \boxed{v_x^2=v_{0x}^2+2a_x\Delta x}. \]

Es especialmente útil cuando el tiempo no es conocido ni constituye la incógnita.
\end{respuesta}

\begin{explicacion}
Partimos de:

\[ v_x=v_{0x}+a_xt. \]

Si:

\[ a_x\neq0, \]

despejamos el tiempo:

\[
t
=
\frac{v_x-v_{0x}}{a_x}.
\]

Ahora utilizamos:

\[
\Delta x
=
\frac{v_{0x}+v_x}{2}\,t.
\]

Sustituimos la expresión del tiempo:

\[
\Delta x
=
\frac{v_{0x}+v_x}{2}
\left(
\frac{v_x-v_{0x}}{a_x}
\right).
\]

Multiplicando:

\[
\Delta x
=
\frac{
(v_{0x}+v_x)(v_x-v_{0x})
}{
2a_x
}.
\]

El numerador es una diferencia de cuadrados:

\[
(v_{0x}+v_x)(v_x-v_{0x})
=
v_x^2-v_{0x}^2.
\]

Entonces:

\[
\Delta x
=
\frac{
v_x^2-v_{0x}^2
}{
2a_x
}.
\]

Multiplicando por \(2a_x\):

\[
2a_x\Delta x
=
v_x^2-v_{0x}^2.
\]

Finalmente:

\[ \boxed{v_x^2=v_{0x}^2+2a_x\Delta x}. \]

Esta ecuación relaciona directamente:

\[
v_{0x},\qquad
v_x,\qquad
a_x,\qquad
\Delta x
\]

sin incluir el tiempo.

Por ello resulta particularmente útil cuando el problema proporciona velocidades, aceleración y desplazamiento, pero no requiere conocer cuánto dura el movimiento.

Si \(a_x=0\), el movimiento deja de ser MRUA con aceleración no nula y se reduce al caso de velocidad constante. En ese caso no debe realizarse el paso algebraico de dividir entre \(a_x\).
\end{explicacion}